% Importado desde: 2026-03-28_Tiempo_Continuo_vs_QW_QCA.tex
\chapter{Derivacion Pedagogica: --- Tiempo continuo vs tiempo discreto: de Hamiltonianos en red a quantum walks y QCA}
\label{ch:tiempo-continuo-vs-qw-qca}

\section{Objetivo}

Despues de SSH, grafeno, Dirac/Weyl en \(3D\) y Nielsen--Ninomiya, el siguiente paso natural del proyecto es decidir \textbf{como se discretiza exactamente la teoria}. Hay dos rutas principales:

\begin{enumerate}[leftmargin=*]
  \item \textbf{espacio discreto y tiempo continuo:} la dinamica se da por un Hamiltoniano \(H\) y una ecuacion de Schrodinger;
  \item \textbf{espacio-tiempo discreto:} la dinamica se da por una evolucion unitaria por pasos, como en \emph{quantum walks} o \emph{quantum cellular automata} (QCA).
\end{enumerate}

La meta de este capitulo es comparar esas dos estrategias paso a paso y mostrar, con apoyo visual, como ambas pueden llevar a una ecuacion efectiva de Dirac, aunque lo hagan de maneras conceptualmente distintas.

\section{Panorama general}

La diferencia central puede resumirse con una sola pregunta:

\begin{displayquote}
\enquote{La variable tiempo es continua desde el principio, o tambien se construye a partir de pasos discretos?}
\end{displayquote}

Si el tiempo es continuo, escribimos
\begin{equation}
\ii \partial_t \psi = H \psi.
\end{equation}

Si el tiempo es discreto, escribimos
\begin{equation}
\psi(t+\Delta t)=U\,\psi(t),
\end{equation}
donde \(U\) es un operador unitario local.

\begin{figure}[ht]
\centering
\begin{tikzpicture}[>=Latex,node distance=1.6cm]
  \node[draw,rounded corners,fill=blue!6,align=center,minimum width=4.1cm,minimum height=1.05cm] (a1) {Red espacial discreta\\tiempo continuo};
  \node[draw,rounded corners,fill=blue!6,align=center,minimum width=4.1cm,minimum height=1.05cm,below=of a1] (a2) {Hamiltoniano \(H\)};
  \node[draw,rounded corners,fill=blue!6,align=center,minimum width=4.1cm,minimum height=1.05cm,below=of a2] (a3) {\(\ii\partial_t\psi=H\psi\)};

  \node[draw,rounded corners,fill=green!8,align=center,minimum width=4.1cm,minimum height=1.05cm,right=4.3cm of a1] (b1) {Espacio-tiempo discreto\\pasos unitarios};
  \node[draw,rounded corners,fill=green!8,align=center,minimum width=4.1cm,minimum height=1.05cm,below=of b1] (b2) {Operador \(U\)};
  \node[draw,rounded corners,fill=green!8,align=center,minimum width=4.1cm,minimum height=1.05cm,below=of b2] (b3) {\(\psi(t+\Delta t)=U\psi(t)\)};

  \draw[->,thick] (a1) -- (a2);
  \draw[->,thick] (a2) -- (a3);
  \draw[->,thick] (b1) -- (b2);
  \draw[->,thick] (b2) -- (b3);

  \node[draw,rounded corners,fill=orange!10,align=center,minimum width=5.4cm,minimum height=1.05cm,below=2.1cm of $(a3)!0.5!(b3)$] (c) {Limite de baja energia y longitudes de onda grandes\\\(\Longrightarrow\) ecuacion efectiva tipo Dirac};
  \draw[->,thick] (a3) -- (c);
  \draw[->,thick] (b3) -- (c);
\end{tikzpicture}
\caption{Las dos rutas principales del proyecto. Ambas pueden conducir a una dinamica efectiva de Dirac, pero no parten del mismo objeto matematico.}
\label{fig:two-routes}
\end{figure}

\section{Estrategia A: espacio discreto y tiempo continuo}

\subsection{La idea basica}

Esta es la estrategia que ya venimos usando en SSH, grafeno y modelos de Dirac/Weyl. El espacio se discretiza en una red, pero el tiempo sigue siendo continuo. La evolucion esta gobernada por
\begin{equation}
\ii \partial_t \psi = H \psi.
\label{eq:sch-time}
\end{equation}

La informacion del sistema vive en un Hamiltoniano \(H\), construido a partir de saltos locales entre sitios, subredes, orbitales o grados de libertad internos.

\subsection{Ejemplo minimo en 1D}

Tomemos el Hamiltoniano ingenuo de Dirac en red de una dimension:
\begin{equation}
H(k)=\frac{1}{a}\sin(ka)\,\sigma_x + m\,\sigma_z.
\label{eq:lattice-dirac-cont}
\end{equation}

Aqui:

\begin{enumerate}[leftmargin=*]
  \item \(a\) es el espaciado de red;
  \item \(k\) es el momento de Bloch;
  \item \(\sigma_x,\sigma_z\) actuan sobre un espacio interno de dos componentes.
\end{enumerate}

El espectro es
\begin{equation}
E(k)=\pm \sqrt{\frac{1}{a^2}\sin^2(ka)+m^2}.
\end{equation}

Si ahora miramos solo excitaciones de baja energia cerca de \(k=0\), usamos
\begin{equation}
\sin(ka)\simeq ka
\qquad (|ka|\ll 1),
\end{equation}
y obtenemos
\begin{equation}
H(k)\simeq k\,\sigma_x + m\,\sigma_z.
\label{eq:dirac-lowk}
\end{equation}

Volviendo a espacio real, \(k\to -\ii \partial_x\), llegamos a
\begin{equation}
\ii \partial_t \psi(x,t)
=
\left(
-\ii \sigma_x \partial_x + m \sigma_z
\right)\psi(x,t),
\label{eq:dirac-11-cont}
\end{equation}
que es la forma canonica de una ecuacion de Dirac en \(1+1\).

\subsection{Que se gana con esta estrategia}

\begin{enumerate}[leftmargin=*]
  \item conecta directamente con materia condensada y con modelos tipo \emph{tight-binding};
  \item permite estudiar bandas, nodos, simetrias de red y linealizaciones;
  \item es la ruta mas natural para SSH, grafeno, semimetales de Dirac y semimetales de Weyl.
\end{enumerate}

\subsection{Que limita esta estrategia}

\begin{enumerate}[leftmargin=*]
  \item el tiempo no forma parte de la estructura discreta;
  \item la relatividad es efectiva, no microscopica;
  \item el problema de Nielsen--Ninomiya aparece enseguida si se intenta interpretar la red como sustrato mas fundamental.
\end{enumerate}

\begin{figure}[ht]
\centering
\begin{tikzpicture}[>=Latex,scale=1.0]
  \foreach \x in {0,1.2,2.4,3.6,4.8}
    \fill[blue!70] (\x,0) circle (2.2pt);
  \draw[very thick] (0,0)--(1.2,0)--(2.4,0)--(3.6,0)--(4.8,0);
  \node[above] at (2.4,0.25) {\small red espacial};

  \draw[->,very thick] (6.0,-0.4) -- (6.0,1.5);
  \node[left] at (6.0,1.5) {\small \(t\)};
  \node[right] at (6.0,0.6) {\small continuo};

  \node[draw,rounded corners,fill=blue!6,minimum width=3.0cm,minimum height=1.0cm] at (9.2,0.5) {\(\ii \partial_t \psi = H\psi\)};
\end{tikzpicture}
\caption{En esta ruta la red es espacial, pero el tiempo sigue siendo un parametro continuo.}
\label{fig:continuous-time-picture}
\end{figure}

\section{Estrategia B: espacio-tiempo discreto con quantum walks}

\subsection{La idea basica}

Ahora damos un paso mucho mas ambicioso. En lugar de escribir una ecuacion diferencial en el tiempo, imponemos una evolucion por pasos:
\begin{equation}
\psi_{n+1}=U\,\psi_n,
\label{eq:unitary-step}
\end{equation}
donde \(n\) cuenta instantes separados por \(\Delta t\).

El objeto fundamental ya no es un Hamiltoniano, sino un operador unitario \(U\). La pregunta pasa a ser:

\begin{displayquote}
\enquote{Que forma debe tener \(U\) para que, cuando \(\Delta t\) y \(a\) son pequenos, emerja una ecuacion de Dirac?}
\end{displayquote}

\subsection{Una caminata cuantica minima}

En una \emph{quantum walk} 1D, el estado tiene dos componentes internas. Podemos llamarlas
\begin{equation}
\lvert \uparrow \rangle,\qquad \lvert \downarrow \rangle,
\end{equation}
y pensar que una de ellas tiende a moverse hacia la derecha y la otra hacia la izquierda.

Un paso elemental se construye con dos piezas:

\begin{enumerate}[leftmargin=*]
  \item una \textbf{rotacion interna} o \emph{coin} \(C(\theta)\);
  \item un \textbf{desplazamiento condicionado} \(S\).
\end{enumerate}

Tomemos
\begin{equation}
C(\theta)=\ee^{-\ii \theta \sigma_y},
\end{equation}
y un operador de desplazamiento tal que
\begin{equation}
S\,\lvert x,\uparrow \rangle=\lvert x+a,\uparrow \rangle,
\qquad
S\,\lvert x,\downarrow \rangle=\lvert x-a,\downarrow \rangle.
\end{equation}

Entonces un paso completo es
\begin{equation}
U = S\,C(\theta).
\label{eq:U-step}
\end{equation}

\subsection{Como actua en espacio de momentos}

En la representacion de Bloch, el corrimiento condicionado se escribe como
\begin{equation}
S(k)=\ee^{-\ii ka \sigma_z}.
\end{equation}

Por tanto,
\begin{equation}
U(k)=\ee^{-\ii ka \sigma_z}\,\ee^{-\ii \theta \sigma_y}.
\label{eq:Uk}
\end{equation}

Ahora viene el paso clave. Si \(ka\) y \(\theta\) son pequenos, expandimos:
\begin{equation}
\ee^{-\ii ka \sigma_z}\simeq I-\ii ka\,\sigma_z,
\qquad
\ee^{-\ii \theta \sigma_y}\simeq I-\ii \theta\,\sigma_y.
\end{equation}

Multiplicando y reteniendo solo primer orden:
\begin{equation}
U(k)\simeq I-\ii\left(ka\,\sigma_z + \theta\,\sigma_y\right).
\label{eq:Uk-expand}
\end{equation}

Pero una evolucion por un paso pequeno tambien puede escribirse como
\begin{equation}
U(k)\simeq \ee^{-\ii \Delta t\, H_{\text{eff}}(k)}
\simeq
I-\ii \Delta t\, H_{\text{eff}}(k).
\end{equation}

Comparando con \eqref{eq:Uk-expand}, identificamos
\begin{equation}
H_{\text{eff}}(k)=\frac{a}{\Delta t}k\,\sigma_z + \frac{\theta}{\Delta t}\sigma_y.
\label{eq:Heff-walk}
\end{equation}

Si definimos
\begin{equation}
c=\frac{a}{\Delta t},
\qquad
m=\frac{\theta}{\Delta t},
\end{equation}
queda
\begin{equation}
H_{\text{eff}}(k)= c\,k\,\sigma_z + m\,\sigma_y.
\end{equation}

Al volver a espacio real, \(k\to -\ii \partial_x\), obtenemos
\begin{equation}
\ii \partial_t \psi(x,t)
=
\left(
-\ii c\,\sigma_z \partial_x + m\,\sigma_y
\right)\psi(x,t),
\label{eq:dirac-qw}
\end{equation}
que es de nuevo una ecuacion tipo Dirac en \(1+1\).

\subsection{La leccion conceptual}

La ecuacion de Dirac no fue puesta a mano. Emergió del producto entre:

\begin{enumerate}[leftmargin=*]
  \item un operador de mezcla interna;
  \item una traslacion condicionada;
  \item el limite de pasos pequenos en espacio y tiempo.
\end{enumerate}

\begin{figure}[ht]
\centering
\begin{tikzpicture}[>=Latex,scale=1.0]
  \foreach \x in {0,1.3,2.6,3.9,5.2}
    \fill[blue!70] (\x,0) circle (2pt);
  \foreach \x in {0.65,1.95,3.25,4.55}
    \fill[green!70!black] (\x,1.0) circle (2pt);

  \draw[->] (-0.5,-0.2) -- (-0.5,1.4) node[above] {\(t\)};
  \node[left] at (-0.5,0) {\small \(n\)};
  \node[left] at (-0.5,1.0) {\small \(n+1\)};

  \draw[->,thick,blue!70!black] (0,0) -- (0.65,1.0);
  \draw[->,thick,blue!70!black] (1.3,0) -- (1.95,1.0);
  \draw[->,thick,blue!70!black] (2.6,0) -- (3.25,1.0);
  \draw[->,thick,red!70!black] (1.3,0) -- (0.65,1.0);
  \draw[->,thick,red!70!black] (2.6,0) -- (1.95,1.0);
  \draw[->,thick,red!70!black] (3.9,0) -- (3.25,1.0);

  \node[below] at (2.6,-0.15) {\small espacio};
  \node[right] at (5.5,0.9) {\small \(\uparrow\) tiende a la derecha};
  \node[right] at (5.5,0.55) {\small \(\downarrow\) tiende a la izquierda};
\end{tikzpicture}
\caption{Esquema minimo de una caminata cuantica 1D. En cada paso temporal discreto, una componente se desplaza hacia la derecha y la otra hacia la izquierda, despues de una mezcla interna.}
\label{fig:quantum-walk}
\end{figure}

\section{De quantum walks a QCA}

\subsection{Que anade un QCA}

Una caminata cuantica trabaja muy bien como sistema de una particula o de pocas componentes. Un \emph{quantum cellular automaton} (QCA) lleva la idea a una formulacion mas estructural:

\begin{enumerate}[leftmargin=*]
  \item el espacio esta discretizado en celdas;
  \item el tiempo avanza por pasos;
  \item la regla de evolucion es local y unitaria;
  \item el mismo esquema puede extenderse a muchos grados de libertad y a formulaciones mas cercanas a un campo discreto.
\end{enumerate}

Conceptualmente, un QCA es la version \enquote{de muchas celdas y muchas componentes} de la intuicion de las quantum walks.

\subsection{Por que esto importa para tu proyecto}

Si lo que se quiere es solo reproducir la fisica efectiva de Dirac en un material, la estrategia Hamiltoniana suele bastar. Pero si lo que se quiere es explorar un \textbf{ansatz de espacio-tiempo discreto}, entonces los QW/QCA tienen una ventaja decisiva:

\begin{displayquote}
\enquote{discretizan la evolucion misma, no solo la geometria espacial.}
\end{displayquote}

\begin{figure}[ht]
\centering
\begin{tikzpicture}[>=Latex,node distance=1.3cm]
  \node[draw,rounded corners,fill=green!8,minimum width=3.8cm,minimum height=0.95cm,align=center] (q1) {celda local};
  \node[draw,rounded corners,fill=green!8,minimum width=3.8cm,minimum height=0.95cm,align=center,below=of q1] (q2) {regla unitaria local};
  \node[draw,rounded corners,fill=green!8,minimum width=3.8cm,minimum height=0.95cm,align=center,below=of q2] (q3) {paso temporal discreto};
  \node[draw,rounded corners,fill=green!8,minimum width=3.8cm,minimum height=0.95cm,align=center,below=of q3] (q4) {limite continuo efectivo};

  \draw[->,thick] (q1) -- (q2);
  \draw[->,thick] (q2) -- (q3);
  \draw[->,thick] (q3) -- (q4);

  \node[right=1.8cm of q2,align=left] {
    \begin{tabular}{l}
      \small localidad \\
      \small causalidad discreta \\
      \small unitariedad exacta
    \end{tabular}
  };
\end{tikzpicture}
\caption{Idea estructural de un QCA. La nocion de causalidad local queda incorporada desde el nivel microscopico.}
\label{fig:qca-logic}
\end{figure}

\section{Comparacion paso a paso}

\subsection{Objeto fundamental}

\begin{enumerate}[leftmargin=*]
  \item \textbf{tiempo continuo:} el objeto fundamental es \(H\);
  \item \textbf{QW/QCA:} el objeto fundamental es \(U\).
\end{enumerate}

\subsection{Ecuacion microscopica}

\begin{enumerate}[leftmargin=*]
  \item \textbf{tiempo continuo:} \(\ii\partial_t\psi = H\psi\);
  \item \textbf{QW/QCA:} \(\psi_{n+1}=U\psi_n\).
\end{enumerate}

\subsection{Como aparece Dirac}

\begin{enumerate}[leftmargin=*]
  \item \textbf{tiempo continuo:} linealizando el Hamiltoniano cerca de un nodo del espacio recíproco;
  \item \textbf{QW/QCA:} expandiendo el operador unitario de un paso para \(ka\ll 1\) y \(\Delta t\) pequeno.
\end{enumerate}

\subsection{Lectura fisica}

\begin{enumerate}[leftmargin=*]
  \item \textbf{tiempo continuo:} ideal para bandas, valles, nodos de Dirac/Weyl, topologia de materiales;
  \item \textbf{QW/QCA:} ideal para explorar una dinamica local discreta con aspiracion mas fundamental.
\end{enumerate}

\subsection{Riesgo conceptual}

\begin{enumerate}[leftmargin=*]
  \item \textbf{tiempo continuo:} puede quedarse en analogia de materia condensada y no tocar la estructura temporal microscopica;
  \item \textbf{QW/QCA:} es mas cercano a un ansatz de espacio-tiempo discreto, pero mas abstracto y menos directamente conectado con materiales especificos.
\end{enumerate}

\section{Cual conviene primero en tu programa}

Dado el orden de tu proyecto, la respuesta mas razonable es:

\begin{enumerate}[leftmargin=*]
  \item primero consolidar la ruta Hamiltoniana, porque conecta naturalmente con SSH, grafeno, Dirac/Weyl y Nielsen--Ninomiya;
  \item despues pasar a QW/QCA, porque ahi la pregunta ya no es solo \enquote{como emerge Dirac}, sino \enquote{como emerge desde una dinamica local discreta de espacio-tiempo}.
\end{enumerate}

En otras palabras:

\begin{displayquote}
\textbf{La ruta de tiempo continuo es el laboratorio conceptual. La ruta QW/QCA es el paso serio hacia la interpretacion de espacio-tiempo discreto.}
\end{displayquote}

\section{Mapa de continuidad del proyecto}

\begin{figure}[ht]
\centering
\begin{tikzpicture}[>=Latex,node distance=1.35cm]
  \node[draw,rounded corners,fill=blue!6,minimum width=3.0cm,minimum height=0.95cm,align=center] (ssh) {SSH\\\(1+1\)};
  \node[draw,rounded corners,fill=blue!6,minimum width=3.0cm,minimum height=0.95cm,align=center,right=of ssh] (gra) {Grafeno\\\(2+1\)};
  \node[draw,rounded corners,fill=blue!6,minimum width=3.0cm,minimum height=0.95cm,align=center,right=of gra] (weyl) {Dirac/Weyl\\\(3D\)};
  \node[draw,rounded corners,fill=orange!10,minimum width=3.3cm,minimum height=0.95cm,align=center,below=of gra] (nn) {Nielsen--Ninomiya};
  \node[draw,rounded corners,fill=green!8,minimum width=3.5cm,minimum height=0.95cm,align=center,below left=1.4cm and 1.3cm of nn] (ham) {red espacial\\tiempo continuo};
  \node[draw,rounded corners,fill=green!8,minimum width=3.5cm,minimum height=0.95cm,align=center,below right=1.4cm and 1.3cm of nn] (qca) {QW / QCA\\tiempo discreto};
  \node[draw,rounded corners,fill=purple!8,minimum width=4.2cm,minimum height=0.95cm,align=center,below=1.5cm of nn] (final) {\(3+1\) discreto\\como objetivo conceptual};

  \draw[->,thick] (ssh) -- (gra);
  \draw[->,thick] (gra) -- (weyl);
  \draw[->,thick] (gra) -- (nn);
  \draw[->,thick] (weyl) -- (nn);
  \draw[->,thick] (nn) -- (ham);
  \draw[->,thick] (nn) -- (qca);
  \draw[->,thick] (ham) -- (final);
  \draw[->,thick] (qca) -- (final);
\end{tikzpicture}
\caption{Ubicacion de este capitulo en la investigacion. Las dos estrategias no son enemigas: funcionan como dos ramas complementarias despues de haber entendido la emergencia de Dirac y el problema del doblamiento.}
\label{fig:project-map}
\end{figure}

\section{Resumen final}

Las conclusiones operativas de este capitulo son:

\begin{enumerate}[leftmargin=*]
  \item \textbf{tiempo continuo} significa: red discreta, Hamiltoniano \(H\), ecuacion diferencial en el tiempo;
  \item \textbf{QW/QCA} significa: red discreta, operador unitario \(U\), evolucion por pasos temporales;
  \item en ambos casos, la ecuacion de Dirac aparece como teoria efectiva de baja energia o de pequenos pasos;
  \item la ruta Hamiltoniana es la mas natural para materiales y para la escalera SSH \(\to\) grafeno \(\to\) Weyl;
  \item la ruta QW/QCA es la mas natural si queremos tomar en serio la idea de un espacio-tiempo discreto subyacente.
\end{enumerate}

Por eso, el proyecto no necesita elegir entre una y otra demasiado pronto. Puede usarlas en secuencia:

\begin{enumerate}[leftmargin=*]
  \item primero, para aprender como emerge Dirac desde una red;
  \item despues, para preguntar si esa red puede reinterpretarse como una dinamica discreta del propio espacio-tiempo.
\end{enumerate}

\begin{thebibliography}{99}

\bibitem{FarrellyShort2014}
T. C. Farrelly y A. J. Short,
\newblock \enquote{Causal fermions in discrete space-time},
\newblock \emph{Physical Review A} \textbf{89}, 012302 (2014).

\bibitem{Farrelly2020}
T. C. Farrelly,
\newblock \enquote{A review of Quantum Cellular Automata},
\newblock \emph{Quantum} \textbf{4}, 368 (2020).

\bibitem{Arrighi2018}
P. Arrighi, G. Di Molfetta, I. M\'arquez-Mart\'in y A. P\'erez,
\newblock \enquote{The Dirac equation as a quantum walk over the honeycomb and triangular lattices},
\newblock \emph{Physical Review A} \textbf{97}, 062111 (2018).

\end{thebibliography}
